% !TeX program = pdflatex
\documentclass[11pt,a4paper]{amsart}

\newcommand{\notelanguage}{finnish}
% Shared preamble for course notes and exercise sets.

\usepackage[T1]{fontenc}
\usepackage{lmodern}
\usepackage[english]{babel}

\usepackage[a4paper,margin=30mm]{geometry}
\usepackage{microtype}
\usepackage{graphicx}
\usepackage{enumitem}

\usepackage{amsmath}
\usepackage{amssymb}
\usepackage{amsthm}
\usepackage{mathtools}
\usepackage{aliascnt}

\usepackage[hidelinks,pdfusetitle]{hyperref}

\numberwithin{equation}{section}
\setlist{itemsep=0.25em,topsep=0.5em}

\theoremstyle{plain}
\newtheorem{theorem}{Theorem}[section]
\newaliascnt{lemma}{theorem}
\newtheorem{lemma}[lemma]{Lemma}
\aliascntresetthe{lemma}
\newaliascnt{proposition}{theorem}
\newtheorem{proposition}[proposition]{Proposition}
\aliascntresetthe{proposition}
\newaliascnt{corollary}{theorem}
\newtheorem{corollary}[corollary]{Corollary}
\aliascntresetthe{corollary}

\theoremstyle{definition}
\newaliascnt{definition}{theorem}
\newtheorem{definition}[definition]{Definition}
\aliascntresetthe{definition}
\newaliascnt{example}{theorem}
\newtheorem{example}[example]{Example}
\aliascntresetthe{example}
\newaliascnt{problem}{theorem}
\newtheorem{problem}[problem]{Problem}
\aliascntresetthe{problem}

\theoremstyle{remark}
\newaliascnt{remark}{theorem}
\newtheorem{remark}[remark]{Remark}
\aliascntresetthe{remark}

\usepackage[nameinlink,noabbrev]{cleveref}

\crefname{theorem}{theorem}{theorems}
\crefname{lemma}{lemma}{lemmas}
\crefname{proposition}{proposition}{propositions}
\crefname{corollary}{corollary}{corollaries}
\crefname{definition}{definition}{definitions}
\crefname{example}{example}{examples}
\crefname{problem}{problem}{problems}
\crefname{remark}{remark}{remarks}

\DeclareMathOperator{\im}{im}
\DeclareMathOperator{\rank}{rank}
\DeclarePairedDelimiter{\abs}{\lvert}{\rvert}
\DeclarePairedDelimiter{\norm}{\lVert}{\rVert}
\DeclarePairedDelimiter{\set}{\lbrace}{\rbrace}

\usepackage{tikz}
\usetikzlibrary{arrows.meta}
\tikzset{
  mapping/.style={x=1cm, y=0.6cm, baseline=(current bounding box.north)},
  set outline/.style={draw=black!35, fill=black!2, line width=0.5pt},
  element/.style={circle, fill=black, inner sep=0pt, outer sep=0pt,
  minimum size=2.6pt},
  map arrow/.style={-{Stealth[length=1.8mm, width=1.2mm]},
  draw=blue!45!black, line width=0.65pt}
}
\DeclareMathOperator{\sign}{sign}

\title{Johdatus yliopistomatematiikkaan\\Harjoituksen 4 ratkaisut}
\date{}
\author{}

\begin{document}

\exerciseheading{Johdatus yliopistomatematiikkaan}{Harjoituksen 4 ratkaisut}

\begin{exercise}
  Olkoot \(f \colon X \to Y\) ja \(B \subseteq Y\). Osoita, että
  \[
    f(f^{-1}(B)) \subseteq B.
  \]
  Anna esimerkki kuvauksesta \(f\) ja joukosta \(B\), jotka osoittavat,
  että tuloksessa osajoukko voi olla aito.
\end{exercise}

\begin{solution}
  Olkoon \(y \in f(f^{-1}(B))\). Funktion kuvan määritelmän mukaan on olemassa
  \(x \in f^{-1}(B)\) siten, että \(y = f(x)\). Funktion alkukuvan
  määritelmän mukaan \(f(x) \in B\), joten \(y = f(x) \in B\).
  Koska \(y\) oli joukon \(f(f^{-1}(B))\)
  mielivaltainen alkio, on voimassa
  \[
    f(f^{-1}(B)) \subseteq B.
  \]
  Olkoon \(f \colon \set{1, 2} \to \set{1, 2}\), \(f(1) = f(2) = 1\).
  Valitaan \(B = \set{2}\). Tällöin
  \[
    f(f^{-1}(B)) = f(\emptyset) = \emptyset \subsetneq B,
  \]
  joten osajoukko voi olla aito.
\end{solution}

\begin{exercise}
  Olkoot \(X = \set{a, b, c, d}\) ja \(Y = \set{1, 2, 3}\). Keksi kuvaus
  \(X \to Y\) tai \(Y \to X\) tai \(X \to X\) tai \(Y \to Y\), joka
\begin{enumerate}[label=\alph*), nosep]
  \item on injektio, mutta ei ole surjektio,
  \item on surjektio, mutta ei ole injektio,
  \item ei ole surjektio eikä injektio,
  \item on bijektio.
\end{enumerate}
Anna vastauksesi nuolikuvioina. Vastauksia ei tarvitse perustella.
\end{exercise}

\begin{solution}
\leavevmode\par
\begin{center}
  \begin{tabular}{@{}c@{\hspace{1.4cm}}c@{}}
    \begin{tikzpicture}[mapping]
      \path[use as bounding box] (-0.55, -3.6) rectangle (2.95, 2.1);
      \draw[set outline] (0, -1) ellipse [x radius=0.43cm, y radius=0.9cm];
      \draw[set outline] (2.4, -1.5) ellipse [x radius=0.43cm, y radius=1.2cm];
    \node at (1.2, 1.8) {a) \(f \colon Y \to X\)};
    \node at (0, 0.8) {\(Y\)};
    \node at (2.4, 0.8) {\(X\)};
    \foreach \i in {1,2,3}
    \node[element, label={[label distance=3pt, inner sep=0pt]left:\(\i\)}]
    (s\i) at (0.14, 1-\i) {};
    \foreach \i/\v in {1/a,2/b,3/c,4/d}
    \node[element, label={[label distance=3pt, inner sep=0pt]right:\(\v\)}]
    (t\i) at (2.26, 1-\i) {};
    \foreach \i in {1,2,3}
    \draw[map arrow] (s\i) -- (t\i);
  \end{tikzpicture}
  &
  \begin{tikzpicture}[mapping]
    \path[use as bounding box] (-0.55, -3.6) rectangle (2.95, 2.1);
    \draw[set outline] (0, -1.5) ellipse [x radius=0.43cm, y radius=1.2cm];
    \draw[set outline] (2.4, -1) ellipse [x radius=0.43cm, y radius=0.9cm];
  \node at (1.2, 1.8) {b) \(g \colon X \to Y\)};
  \node at (0, 0.8) {\(X\)};
  \node at (2.4, 0.8) {\(Y\)};
  \foreach \i/\v in {1/a,2/b,3/c,4/d}
  \node[element, label={[label distance=3pt, inner sep=0pt]left:\(\v\)}]
  (s\i) at (0.14, 1-\i) {};
  \foreach \i in {1,2,3}
  \node[element, label={[label distance=3pt, inner sep=0pt]right:\(\i\)}]
  (t\i) at (2.26, 1-\i) {};
  \foreach \i in {1,2,3}
  \draw[map arrow] (s\i) -- (t\i);
  \draw[map arrow] (s4) -- (t3);
\end{tikzpicture}
\\[0.5em]
\begin{tikzpicture}[mapping]
  \path[use as bounding box] (-0.55, -3.6) rectangle (2.95, 2.1);
  \draw[set outline] (0, -1) ellipse [x radius=0.43cm, y radius=0.9cm];
  \draw[set outline] (2.4, -1) ellipse [x radius=0.43cm, y radius=0.9cm];
\node at (1.2, 1.8) {c) \(h \colon Y \to Y\)};
\node at (0, 0.8) {\(Y\)};
\node at (2.4, 0.8) {\(Y\)};
\foreach \i in {1,2,3} {
  \node[element, label={[label distance=3pt, inner sep=0pt]left:\(\i\)}]
  (s\i) at (0.14, 1-\i) {};
  \node[element, label={[label distance=3pt, inner sep=0pt]right:\(\i\)}]
  (t\i) at (2.26, 1-\i) {};
}
\foreach \i in {1,2,3}
\draw[map arrow] (s\i) -- (t1);
\end{tikzpicture}
&
\begin{tikzpicture}[mapping]
\path[use as bounding box] (-0.55, -3.6) rectangle (2.95, 2.1);
\draw[set outline] (0, -1) ellipse [x radius=0.43cm, y radius=0.9cm];
\draw[set outline] (2.4, -1) ellipse [x radius=0.43cm, y radius=0.9cm];
\node at (1.2, 1.8) {d) \(k \colon Y \to Y\)};
\node at (0, 0.8) {\(Y\)};
\node at (2.4, 0.8) {\(Y\)};
\foreach \i in {1,2,3} {
\node[element, label={[label distance=3pt, inner sep=0pt]left:\(\i\)}]
(s\i) at (0.14, 1-\i) {};
\node[element, label={[label distance=3pt, inner sep=0pt]right:\(\i\)}]
(t\i) at (2.26, 1-\i) {};
\draw[map arrow] (s\i) -- (t\i);
}
\end{tikzpicture}
\end{tabular}
\end{center}
\end{solution}

\newpage
\begin{exercise}
Tutki, ovatko seuraavat kuvaukset injektioita, surjektioita ja/tai
bijektioita:
\begin{enumerate}[label=\alph*), nosep]
\item \(f \colon \NN \to \NN\), \(f(x) = 3x\),
\item \(g \colon \ZZ \to \NN\), \(g(x) = \abs{x}\),
\item \(h \colon \RR \to \RR\), \(h(x) = -x + 1\),
\end{enumerate}
Huomioi erityisesti kuvausten lähtö- ja maalijoukot!
\end{exercise}

\begin{solution}
\leavevmode\par
\begin{enumerate}[label=\alph*), nosep]
\item Olkoot \(x_1, x_2 \in \NN\) ja oletetaan, että \(f(x_1) = f(x_2)\).
Tällöin
\[
f(x_1) = f(x_2) \implies 3x_1 = 3x_2 \implies x_1 = x_2,
\]
joten \(f\) on injektio. Kuvaus \(f\) ei ole surjektio, sillä
\(2 \in \NN\), mutta yhtälöstä \(f(x) = 2\) seuraa
\(x = \frac{2}{3} \notin \NN\). Täten \(f\) ei ole bijektio.
\item Kuvaus \(g\) ei ole injektio, sillä \(-1, 1 \in \ZZ\) ja
\(-1 \neq 1\), mutta
\[
g(-1) = \abs{-1} = 1 = \abs{1} = g(1).
\]
Olkoon \(y \in \NN\) mielivaltainen. Valitaan \(x = y \in \ZZ\).
Koska \(y \geq 0\), saadaan
\[
g(x) = \abs{y} = y,
\]
joten \(g\) on surjektio. Koska \(g\) ei ole injektio, se ei ole bijektio.
\item Olkoot \(x_1, x_2 \in \RR\) ja oletetaan, että \(h(x_1) = h(x_2)\).
Tällöin
\[
h(x_1) = h(x_2) \implies -x_1 + 1 = -x_2 + 1 \implies x_1 = x_2,
\]
joten \(h\) on injektio. Olkoon \(y \in \RR\) mielivaltainen.
Valitaan \(x = 1-y \in \RR\). Tällöin
\[
h(x) = -(1-y) + 1 = y,
\]
joten \(h\) on surjektio. Täten \(h\) on bijektio.
\end{enumerate}
\end{solution}

\begin{exercise}
Olkoot \(f, g \colon \RR \to \RR\), \(f(x) = x + 2\) ja
\(g(x) = \abs{x - 1}\).
\begin{enumerate}[label=\alph*), nosep]
\item Määritä funkion \(f\) käänteisfunktio
\(f^{-1} \colon \RR \to \RR\), jos se on olemassa.
\item Määritä funkion \(g\) käänteisfunktio
\(g^{-1} \colon \RR \to \RR\), jos se on olemassa.
\item Määritä yhdistetyt funktiot \(f \circ g\) ja \(g \circ f\).
Onko \(f \circ g = g \circ f\)?
\end{enumerate}
\end{exercise}

\begin{exercise}
Todista suljettujen reaalilukuvälien yhtä mahtavuus osoittamalla, että
suljetut välit \([0, a]\) ja \([0, b]\) ovat yhtä mahtavia, kun
\(a, b > 0\). \emph{Vihje:} Tutki kuvausta
\(f \colon [0, a] \to [0, b]\), \(f(x) = bx/a\).
\end{exercise}

\begin{exercise}
Onko kahden numeroituvan joukon karteesinen tulo numeroituva? Perustele.
\end{exercise}

\end{document}
